% Extra LaTeX loaded into the preamble by build-pdf.sh. The document class is
% KOMA-Script's scrartcl, so headings, captions, the TOC and running heads are
% styled through KOMA's own interfaces rather than extra packages.
\makeatletter

% Tables are rendered as raw LaTeX by pdf/tables.lua, so pandoc no longer sees
% them and won't load these itself
\usepackage{longtable,booktabs,array,calc,colortbl,ragged2e,caption}
\usepackage{float,framed}

% ---------------------------------------------------------------------------
% Palette
% ---------------------------------------------------------------------------
\definecolor{brand}{HTML}{1F365B}      % navy: title page, headings, table heads
\definecolor{accent}{HTML}{2C6DB5}     % section numbers, links, bullets
\definecolor{accentlight}{HTML}{8DB8E8}% title page highlight
\definecolor{muted}{HTML}{5B6575}      % running heads, captions
\definecolor{hairline}{HTML}{D6DCE5}   % rules and figure frames
\definecolor{tablerow}{HTML}{F2F5F9}   % zebra stripe

% hyperref is loaded after this file, so configure it once the document starts
\AtBeginDocument{\hypersetup{linkcolor=brand, urlcolor=accent,
  citecolor=accent, pdfdisplaydoctitle=true}}

% ---------------------------------------------------------------------------
% Headings: navy, accent-coloured numbers, hairline under sections
% ---------------------------------------------------------------------------
\setkomafont{disposition}{\color{brand}\normalfont\bfseries\hypersetup{urlcolor=brand}}
\setkomafont{section}{\LARGE}
\setkomafont{subsection}{\Large}
\setkomafont{subsubsection}{\large}
\renewcommand*{\sectionformat}{\textcolor{accent}{\thesection}\enskip}
\renewcommand*{\subsectionformat}{\textcolor{accent}{\thesubsection}\enskip}
\renewcommand*{\subsubsectionformat}{\textcolor{accent}{\thesubsubsection}\enskip}
\RedeclareSectionCommand[beforeskip=-1.6\baselineskip, afterskip=.6\baselineskip]{section}
\RedeclareSectionCommand[beforeskip=-1.1\baselineskip, afterskip=.35\baselineskip]{subsection}
\RedeclareSectionCommand[beforeskip=-.9\baselineskip, afterskip=.25\baselineskip]{subsubsection}
\renewcommand*{\sectionlinesformat}[4]{%
  \Ifstr{#1}{section}{%
    \parbox[t]{\linewidth}{\raggedsection\@hangfrom{\hskip #2#3}{#4}\par
      \kern 5pt{\color{hairline}\hrule height .8pt}}%
  }{\@hangfrom{\hskip #2#3}{#4}}%
}

% ---------------------------------------------------------------------------
% Title page
% ---------------------------------------------------------------------------
\renewcommand{\maketitle}{%
  \begin{titlepage}
    \pagecolor{brand}\color{white}\normalfont\setstretch{1}\raggedright
    {\small\addfontfeature{LetterSpace=18}%
      \MakeUppercase{\ifx\@subtitle\@empty Report\else\@subtitle\fi}\par}
    \vspace*{\fill}
    {\color{accentlight}\rule{2.5cm}{3pt}\par}
    \vspace{1.6em}
    {\fontsize{34}{41}\selectfont\bfseries\@title\par}
    \vspace{2.4em}
    {\LARGE\@author\par}
    \vspace*{1.6\fill}
    {\color{white!30!brand}\rule{\linewidth}{.5pt}\par}
    \vspace{.8em}
    {\small\color{white!65!brand}%
      Updated \ifx\@date\@empty\today\else\@date\fi\par}
  \end{titlepage}
  \nopagecolor
  \pagenumbering{roman}%
}

% ---------------------------------------------------------------------------
% Front matter: contents, figures and tables flow together on roman-numbered
% pages, then the body starts on a fresh page at 1. build-pdf.sh only passes
% --toc; the lists of figures and tables are added here.
% ---------------------------------------------------------------------------
% The figure and table counts are saved to the .aux file at the end of each
% run (pandoc reruns LaTeX for the TOC anyway), so empty lists are skipped
\AtEndDocument{\immediate\write\@auxout{%
  \gdef\string\figurecount{\the\value{figure}}%
  \gdef\string\tablecount{\the\value{table}}}}
\def\figurecount{1}\def\tablecount{1}
\let\oldtableofcontents\tableofcontents
\renewcommand{\tableofcontents}{%
  \pagestyle{plain.scrheadings}%
  \oldtableofcontents
  \ifnum\figurecount>0 \bigskip\listoffigures\fi
  \ifnum\tablecount>0 \bigskip\listoftables\fi
  \clearpage
  \pagenumbering{arabic}%
  \pagestyle{scrheadings}%
}
\DeclareTOCStyleEntry[entryformat=\normalfont\bfseries\color{brand},
  pagenumberformat=\normalfont\bfseries\color{brand},
  beforeskip=.9em plus 1pt]{tocline}{section}
\DeclareTOCStyleEntries[entryformat=\normalfont,
  pagenumberformat=\normalfont\color{muted}]{tocline}{subsection,subsubsection}
\DeclareTOCStyleEntries[entryformat=\normalfont,
  pagenumberformat=\normalfont\color{muted}]{tocline}{figure,table}

% Accent-coloured list bullets to match the tables
\renewcommand{\labelitemi}{\textcolor{accent}{\textbullet}}

% ---------------------------------------------------------------------------
% Running heads: title and current section above, author and page below
% ---------------------------------------------------------------------------
\usepackage[headsepline=.4pt, footsepline=.4pt]{scrlayer-scrpage}
\clearpairofpagestyles
\automark[section]{section}
\setkomafont{pageheadfoot}{\normalfont\footnotesize\color{muted}}
\setkomafont{pagenumber}{\normalfont\footnotesize\bfseries\color{brand}}
\setkomafont{headsepline}{\color{hairline}}
\setkomafont{footsepline}{\color{hairline}}
\ihead{\@title}
\ohead{\headmark}
\ifoot{\@author}
\ofoot[\pagemark]{\pagemark}
\pagestyle{scrheadings}

% ---------------------------------------------------------------------------
% Captions: small with a bold navy label
% ---------------------------------------------------------------------------
\captionsetup{font=small, labelfont={bf,color=brand},
  textfont={color=muted}, labelsep=quad, skip=8pt}
\captionsetup[table]{position=top, justification=raggedright, singlelinecheck=false,
  skip=8pt}

% ---------------------------------------------------------------------------
% Figures stay where they appear in the markdown, framed with a hairline
% ---------------------------------------------------------------------------
\floatplacement{figure}{H}
\AtBeginEnvironment{figure}{%
  \let\plainincludegraphics\includegraphics
  \renewcommand{\includegraphics}[2][]{%
    {\setlength{\fboxsep}{0pt}\setlength{\fboxrule}{.6pt}%
     \fcolorbox{hairline}{white}{\plainincludegraphics[#1]{#2}}}}%
}

% ---------------------------------------------------------------------------
% Tables (rendered by pdf/tables.lua): navy header band, zebra rows, small
% text, a coloured bottom rule and real hanging bullets inside cells
% ---------------------------------------------------------------------------
\newcommand{\styledtable}{%
  \normalfont\small\setstretch{1.15}%
  \setlength{\parskip}{0pt}%
  \setlength{\tabcolsep}{7pt}%
  \setlength{\extrarowheight}{5pt}%
  \renewcommand{\arraystretch}{1}%
  \arrayrulecolor{brand}%
}
\newenvironment{longtablewrap}{%
  \par\medskip\styledtable
  \setlength{\LTpre}{0pt}\setlength{\LTpost}{0pt}%
}{\par\bigskip}
\newcommand{\tablestrut}{\rule[-7pt]{0pt}{0pt}}
\newcommand{\tablehead}{\leavevmode\color{white}\bfseries}
% XeTeX lets the coloured longtable caption leak into the rows, so every body
% cell resets the colour
\newcommand{\tablebody}{\leavevmode\normalcolor}
% The invisible strut line gives the cell a first baseline so the first item
% lines up with text in neighbouring cells
\newenvironment{tablelist}{%
  \strut\par\vspace*{-\baselineskip}%
  \begin{list}{\textcolor{accent}{\textbullet}}{%
    \setlength{\leftmargin}{1em}\setlength{\labelwidth}{.6em}%
    \setlength{\labelsep}{.4em}\setlength{\itemsep}{2pt}%
    \setlength{\parsep}{0pt}\setlength{\topsep}{0pt}%
    \setlength{\partopsep}{0pt}}%
}{\end{list}}

% ---------------------------------------------------------------------------
% Callouts: GitHub-style alerts (pdf/alerts.lua) and plain block quotes as a
% tinted box with a coloured left bar
% ---------------------------------------------------------------------------
\definecolor{alertnote}{HTML}{0969DA}
\definecolor{alerttip}{HTML}{1A7F37}
\definecolor{alertimportant}{HTML}{8250DF}
\definecolor{alertwarning}{HTML}{9A6700}
\definecolor{alertcaution}{HTML}{CF222E}
\newenvironment{alert}[2]{%
  \def\FrameCommand{{\color{#1}\vrule width 3pt}%
    \fboxsep=10pt\colorbox{#1!7}}%
  \MakeFramed{\advance\hsize-\width\FrameRestore}%
  \setstretch{1.3}%
  \noindent{\normalfont\small\bfseries\color{#1}%
    \addfontfeature{LetterSpace=6}\MakeUppercase{#2}}\par\smallskip
}{\endMakeFramed}
\renewenvironment{quote}{%
  \def\FrameCommand{{\color{hairline}\vrule width 3pt}\hspace{10pt}}%
  \MakeFramed{\advance\hsize-\width\FrameRestore}\color{muted}\itshape
}{\endMakeFramed}

% Code blocks: light grey panel (pandoc only defines Shaded when the document
% contains a code block)
\definecolor{codebg}{HTML}{F4F6F9}
\AtBeginDocument{\@ifundefined{Shaded}{}{%
  \renewenvironment{Shaded}{%
    \def\FrameCommand{\fboxsep=8pt\colorbox{codebg}}%
    \MakeFramed{\advance\hsize-\width\FrameRestore}%
    \small\setstretch{1.15}%
  }{\endMakeFramed}}}

\makeatother
